\chapter{Resultados Obtidos}
\label{cap:resultados}

O desenvolvimento do GranaAI reuniu três resultados principais: o registro de movimentações financeiras por mensagens, a integração desses registros à plataforma web e a disponibilização de recursos de organização e acompanhamento financeiro. A análise a seguir considera as operações e as interfaces apresentadas no Capítulo~\ref{cap:solucao}, relacionando-as às necessidades que orientaram o projeto.

No registro pelo WhatsApp, o resultado pode ser observado no exemplo da Figura~\ref{fig:md-15}. A mensagem ``Pizza 45 reais'' deu origem a um resumo com o valor da despesa, sua descrição e a categoria Alimentação. Após a confirmação, o assistente informou a conclusão do registro na própria conversa. O exemplo apresenta o fluxo desde o envio de uma informação em linguagem cotidiana até a confirmação da movimentação, reunindo em uma interação tarefas que também podem ser realizadas por formulário na plataforma web.

A confirmação permite conferir o conteúdo interpretado antes de concluir a operação. Esse procedimento é relevante porque o registro depende da identificação das informações presentes na mensagem, que pode variar conforme a forma de escrita. No fluxo apresentado, a conferência antecede a gravação da despesa e mantém o usuário envolvido na definição dos dados que serão armazenados.

A integração com a plataforma web complementa esse processo. Durante o desenvolvimento, foi verificado o fluxo de registro de movimentações pelo WhatsApp e sua apresentação no dashboard. Dessa maneira, o canal de entrada da informação e o ambiente de acompanhamento passam a utilizar os mesmos registros financeiros. Esse resultado corresponde à proposta de combinar a comunicação por mensagens com uma interface voltada à consulta e ao gerenciamento dos dados.

Na plataforma web, o painel e a listagem de transações oferecem níveis distintos de consulta. Os indicadores e gráficos reúnem informações consolidadas, enquanto a listagem permite acompanhar as movimentações individualmente. O formulário de nova transação acrescenta uma alternativa de entrada de dados ao registro conversacional (Figuras~\ref{fig:md-16}, \ref{fig:prints-transacoes} e~\ref{fig:prints-nova-transacao}). Assim, as duas interfaces atendem a atividades complementares dentro do sistema.

Quanto à organização financeira, foram implementados recursos de categorias e limites mensais, metas e assinaturas. As categorias permitem agrupar movimentações e consultar os valores associados aos limites definidos. As metas acrescentam o acompanhamento de objetivos por meio dos valores pretendidos e acumulados. Já as assinaturas reúnem compromissos recorrentes, seus valores e as datas previstas para cobrança. Os exemplos apresentados nas Figuras~\ref{fig:md-17} a~\ref{fig:prints-nova-assinatura} mostram como esses recursos foram incorporados à plataforma, ampliando seu escopo para além da listagem de receitas e despesas.

Parte dessas informações também pode ser consultada na conversa. Na Figura~\ref{fig:wa-categorias-metas}, o assistente apresenta categorias, limites e progresso das metas. Na Figura~\ref{fig:wa-assinaturas}, a solicitação ``Minhas assinaturas'' retorna os serviços acompanhados, os valores e os dias previstos para cobrança. Esses exemplos mostram a aplicação da interface conversacional tanto ao registro quanto à consulta de informações já disponíveis no sistema.

Os relatórios acrescentam uma visão do período financeiro por meio de receitas, despesas e saldo, acompanhados de gráficos e informações analíticas (Figura~\ref{fig:md-19}). A projeção apresentada na Figura~\ref{fig:prints-projecao} estende essa visualização para períodos posteriores. Trata-se de uma estimativa produzida a partir dos dados utilizados pelo sistema, cuja apresentação não representa garantia de saldo futuro. A área de perfil, por sua vez, reúne os dados da conta e as preferências do usuário, complementando os recursos de gerenciamento da aplicação (Figura~\ref{fig:prints-perfil}).

Em relação à pesquisa exploratória, o registro por mensagens corresponde à busca por formas práticas e familiares de anotar informações. Os recursos de categorias, metas, assinaturas e relatórios contemplam outras necessidades de organização identificadas no levantamento. As 35 respostas contribuíram, portanto, para transformar as dificuldades investigadas em requisitos e funcionalidades. Como a coleta antecedeu o desenvolvimento da aplicação atual, essa relação expressa a orientação adotada no projeto.

Os resultados apresentados neste capítulo abrangem as funcionalidades implementadas e as operações documentadas durante o desenvolvimento. Os exemplos de registro, consulta e acompanhamento financeiro mostram a aplicação dos requisitos na solução desenvolvida, desde o envio de mensagens pelo WhatsApp até a apresentação das informações na plataforma web.

O principal resultado técnico é, assim, a integração do registro conversacional aos recursos de organização e acompanhamento da plataforma web. O WhatsApp oferece um canal de entrada e consulta, enquanto a plataforma reúne os dados para gerenciamento e visualização. Essa combinação concretiza a proposta de manter as duas formas de interação vinculadas às mesmas informações financeiras.

\section{Trabalhos Futuros}
\label{sec:trabalhos-futuros}

Como continuidade do trabalho, está planejada a integração do GranaAI ao Open Finance, com o objetivo de permitir a comunicação entre as instituições financeiras e a plataforma, mediante autorização do usuário. Essa integração busca reduzir uma etapa entre a realização de uma despesa e seu registro no sistema, por meio da importação das movimentações disponibilizadas pelas instituições, diminuindo a necessidade de informar cada operação por mensagem ou formulário. Sua implementação deverá contemplar a conciliação dos dados importados com os registros existentes, para evitar duplicidades.

Também está prevista a implementação de planos de assinatura para a monetização da plataforma. Essa etapa deverá incluir a definição dos recursos oferecidos em cada plano, das condições de acesso e dos mecanismos de cobrança, buscando viabilizar a manutenção e a continuidade do serviço.

Outra frente de trabalho é a ampliação dos testes do sistema, incluindo os fluxos de registro e consulta, a integração entre o WhatsApp e a plataforma web e o tratamento de falhas. Para a interpretação das mensagens, deverão ser considerados casos com diferentes formas de descrever valores, categorias e movimentações financeiras. Também se prevê a realização de testes de usabilidade com usuários, para identificar dificuldades no registro, na consulta e na navegação entre os recursos da aplicação.

Em relação aos modelos de linguagem, pretende-se investigar o ajuste fino, ou \textit{fine-tuning}, para adaptá-los às mensagens utilizadas no contexto do controle financeiro pessoal. O objetivo é melhorar a precisão na identificação das intenções e das informações enviadas pelos usuários e reduzir o tempo de resposta. A adoção dessa técnica dependerá da compatibilidade com os modelos e a infraestrutura escolhidos, e seus efeitos deverão ser avaliados por meio da comparação das taxas de acerto e dos tempos de processamento antes e após o ajuste.

Por fim, propõe-se acompanhar a utilização do GranaAI no cotidiano ao longo de um período definido, para investigar sua contribuição para a manutenção de uma rotina de controle financeiro. Essas avaliações ainda não foram realizadas e deverão considerar critérios e procedimentos próprios, distintos do levantamento exploratório que orientou o desenvolvimento.
